\begin{derivation}[轴向变换中的费米子测度与反常]
在欧几里得时空中取规范表示 $R$，并令
\begin{equation*}
 D_\mu=\partial_\mu+\ii g\mathcal A_\mu^aT^a,
 \qquad
 \Tr_R(T^aT^b)=T(R)\delta^{ab}.
\end{equation*}
用厄米算符 $\mathcal H\equiv\ii\slashed D_E$ 的正交本征函数展开 Grassmann 场：
\begin{equation*}
 \mathcal H\phi_n=\lambda_n\phi_n,
 \qquad
 \psi=\sum_na_n\phi_n,
 \qquad
 \bar\psi=\sum_n\bar b_n\phi_n^\dagger.
\end{equation*}
为了定义量子测度沿轴向方向的变化，取一参数局域轴向轨道
\begin{equation*}
 \psi_\tau=\ee^{\ii\tau\alpha(x)\gamma^5}\psi,
 \qquad
 \bar\psi_\tau=\bar\psi\,\ee^{\ii\tau\alpha(x)\gamma^5},
 \qquad \tau\in\mathbb R.
\end{equation*}
在本征函数系数空间中，$\psi$ 的变换矩阵为
\begin{equation*}
 C_{mn}(\tau)
 \equiv\int\dd^4x\,
 \phi_m^\dagger(x)\ee^{\ii\tau\alpha(x)\gamma^5}\phi_n(x),
 \qquad C(0)=I,
\end{equation*}
其在恒等元处的切向量严格为
\begin{equation*}
 \left.\frac{\dd C_{mn}}{\dd\tau}\right|_{\tau=0}
 =\ii\int\dd^4x\,
 \alpha(x)\phi_m^\dagger\gamma^5\phi_n.
\end{equation*}
Grassmann 测度按行列式的逆变换；$\psi$ 与 $\bar\psi$ 两份测度相乘。令 $J(\tau)$ 为总 Jacobian，则
\begin{align*}
 \left.\frac{\dd}{\dd\tau}\ln J(\tau)\right|_{\tau=0}
 &=-2\Tr\!\left[\left.\frac{\dd C}{\dd\tau}\right|_{0}\right]\\
 &=-2\ii\int\dd^4x\,\alpha(x)
 \sum_n\phi_n^\dagger(x)\gamma^5\phi_n(x).
\end{align*}
Ward 恒等式只需要这一沿群轨道的切向变化；随后对迹作规范协变正则化。
模和紫外发散，必须用保持规范协变性的同一算符谱正则化：
\begin{equation*}
 \mathscr D_R(x)
 \equiv\lim_{M\to\infty}
 \tr_{\rm spin,R}\!\left[
 \gamma^5\langle x|\ee^{-\mathcal H^2/M^2}|x\rangle
 \right].
\end{equation*}
把欧氏约定固定到每一个负号。令 $\mu=0,1,2,3$、
$\epsilon_E^{0123}=+1$，采用
\begin{align*}
\{\gamma_E^\mu,\gamma_E^\nu\}&=2\delta^{\mu\nu}I,\qquad
(\gamma_E^\mu)^\dagger=\gamma_E^\mu,\\
\gamma_E^5&=-\gamma_E^0\gamma_E^1\gamma_E^2\gamma_E^3,\qquad
\sigma_E^{\mu\nu}=\frac1{2\ii}[\gamma_E^\mu,\gamma_E^\nu].
\end{align*}
欧氏手征定向按后文的指数 $n_L-n_R$ 选择，即欧氏 $\gamma_E^5$ 的正本征空间
记为 $L$。它的正负标签须随 Wick 旋转一起匹配，不能直接代入闵氏
$P_L=(1-\gamma_M^5)/2$ 的投影符号。
欧氏 $\gamma$ 矩阵厄米，$D_\mu$ 与 $\slashed D_E$ 反厄米；
$M$ 是逆长度量纲的谱截断尺度。协变导数的曲率满足
$[D_\mu,D_\nu]=\ii g\mathcal F_{\mu\nu}^aT^a$，所以
\begin{align*}
\slashed D_E^2
&=D^2+\frac14[\gamma_E^\mu,\gamma_E^\nu][D_\mu,D_\nu]
=D^2-\frac g2\sigma_E^{\mu\nu}\mathcal F_{\mu\nu}^aT^a,\\
\mathcal H^2&=-(D^2+E),\qquad
E=-\frac g2\sigma_E^{\mu\nu}\mathcal F_{\mu\nu}^aT^a.
\end{align*}
Pauli 曲率项由 $\sigma_E^{\mu\nu}$ 的定义固定，手征迹则还依赖 $\gamma_E^5$ 的定向。若把 $\gamma_E^5$ 变号，手征迹及指数的左右手标签随之变号，而 Pauli 项不变。

对光滑规范背景，取无边界紧致调节，或在离边界固定距离的局域区域计算。
令 $t=M^{-2}$、$\Omega_{\mu\nu}=[D_\mu,D_\nu]$。Laplace 型算符的同点热核为
\begin{equation*}
\langle x|\ee^{-t\mathcal H^2}|x\rangle
=\frac1{(4\pi t)^2}\left[
I+tE+t^2\left(\frac12E^2+\frac16D^2E
+\frac1{12}\Omega_{\mu\nu}\Omega_{\mu\nu}\right)+\order(t^3)\right].
\end{equation*}
这里调用光滑背景的局域热核展开；其系数和适用边界条件参见
\href{https://arxiv.org/abs/hep-th/0306138}{Vassilevich 的热核展开综述}第 4、7 节。
该数学输入规定了曲率项、协变导数项与短时余项；以下直接对它们逐项取手征迹。
\begin{align*}
\tr_{\rm spin}(\gamma_E^5 I)&=0,&
\tr_{\rm spin}(\gamma_E^5 E)&=0,\\
\tr_{\rm spin}(\gamma_E^5D^2E)&=0,&
\tr_{\rm spin}(\gamma_E^5\Omega_{\mu\nu}\Omega_{\mu\nu})&=0.
\end{align*}
前两个等式因为不足四个不同的 $\gamma$ 矩阵；导数不会增加 $E$ 的自旋矩阵数；
$\Omega$ 则在自旋空间正比于单位矩阵。只有 $E^2/2$ 的迹留下四个不同指标。

由 $\sigma_E^{\mu\nu}=-\ii\gamma_E^\mu\gamma_E^\nu$（$\mu\ne\nu$）以及
$\gamma_E^5=-\gamma_E^0\gamma_E^1\gamma_E^2\gamma_E^3$，逐项反对易给
\begin{equation*}
\tr_{\rm spin}(\gamma_E^5\sigma_E^{\mu\nu}\sigma_E^{\rho\sigma})
=4\epsilon_E^{\mu\nu\rho\sigma}.
\end{equation*}
指标重复时迹为零；四指标不同时只由排列奇偶性决定符号。于是
\begin{align*}
\mathscr D_R(x)
&=\frac1{16\pi^2}\frac12\frac{g^2}{4}
\left(4\epsilon_E^{\mu\nu\rho\sigma}\right)
\Tr_R(T^aT^b)\mathcal F_{\mu\nu}^a\mathcal F_{\rho\sigma}^b\\
&=\frac{g^2T(R)}{32\pi^2}
\epsilon_E^{\mu\nu\rho\sigma}\mathcal F_{\mu\nu}^a\mathcal F_{\rho\sigma}^a\\
&=\frac{g^2T(R)}{16\pi^2}
\mathcal F_{\mu\nu}^a{}^\star\!\mathcal F_E^{a\mu\nu}.
\end{align*}
在固定光滑背景下，热核中余项乘上 $t^{-2}$ 后为 $\order(t)=\order(M^{-2})$，
所以去除截断的极限只保留上式的有限曲率二次项。
按上述手征与定向约定继续延拓到正文的闵氏指标，局域轴流 Ward 系数为
$g^2T(R)/(16\pi^2)$；场强、电场分量和四维积分测度必须沿同一 Wick 路径一起延拓。
\begin{equation*}
\mathscr D_R(x)
=\frac{g^2T(R)}{16\pi^2}
\mathcal F_{\mu\nu}^a{}^\star\!\mathcal F^{a\mu\nu}.
\end{equation*}
把这个极限代回测度的切向变化
\begin{equation*}
 \left.\frac{\dd}{\dd\tau}\ln J(\tau)\right|_{0}
 =-2\ii\int\dd^4x\,\alpha(x)\mathscr D_R(x),
\end{equation*}
再与质量项沿同一轴向轨道的经典变化相加，Ward 恒等式为
\begin{equation*}
 \partial_\mu J_5^\mu
 =\frac{2\ii mc}{\hbar}\bar\psi\gamma^5\psi
 +\frac{g^2\,2T(R)}{16\pi^2}
 \mathcal F_{\mu\nu}^a{}^\star\!\mathcal F^{a\mu\nu}.
\end{equation*}
QCD 基本表示满足 $T_F=1/2$；对 $N_f$ 个 Dirac 味求和即得到正文
\eqnrefs{eq:qcd-singlet-axial-anomaly-dp32}。两个 Grassmann 测度给出的因子 $2$
与基本表示的 $T_F=1/2$ 恰好相消。
对单位电荷阿贝尔表示则 $T(R)=1$，系数是 $g_{\rm em}^2/(8\pi^2)$。
再以 $F^{\rm geo}=\sqrt{\varepsilon_0c/\hbar}\,F$ 和
$g_{\rm em}=e/\sqrt{\varepsilon_0\hbar c}$ 恢复 SI 单位，
\begin{equation*}
\frac{g_{\rm em}^2}{8\pi^2}F^{\rm geo}{}^\star F^{\rm geo}
=\frac{e^2}{8\pi^2\hbar^2}F{}^\star F,
\end{equation*}
与正文\eqnrefs{eq:abj-anomaly-qed-dp11} 完全一致。反常来自有限的曲率二次热核系数；
光滑背景与规范协变调节是这一计算的假设，不能把它替换成不保持矢量规范对称的截断。
\end{derivation}

\begin{derivation}[电弱拓扑跃迁与重子轻子数选择定则]
先把拓扑信息与费米子路径积分连接起来。对一个处在 $SU(2)$ 表示 $R$ 的左手 Weyl 场，欧几里得 Dirac 算符的指数为
\begin{equation*}
 \operatorname{ind}\slashed D_E
 \equiv n_L-n_R=2T(R)\nu.
\end{equation*}
标准模型弱二重态属于 $R=\mathbf2$，$T(\mathbf2)=1/2$。在自对偶瞬子的半经典背景中，相反手征满足零模消失条件；当 $\nu>0$ 时，每个左手二重态因而有 $\nu$ 个左手零模而没有右手零模。一般光滑拓扑背景的指数只固定两者之差，允许额外的成对零模：
\begin{equation*}
 n_L-n_R=\nu.
\end{equation*}

把一个左手二重态沿 Dirac 本征函数展开，显式分离零模系数，
\begin{equation*}
 \psi_L(x)=\sum_{r=1}^{\nu}a_r\phi_{0r}(x)
 +\sum_{\lambda_n\neq0}a_n\phi_n(x).
\end{equation*}
费米子二次作用量只含非零本征值，因而测度中零模方向留下独立 Grassmann 积分
\begin{equation*}
 \int\mathcal D\psi_L\,\ee^{-S_F}
 =\mathcal N_{\rm nzm}
 \left(\prod_{r=1}^{\nu}\int\dd a_r\right),
 \qquad
 \mathcal N_{\rm nzm}\equiv
 \mathcal N_{\rm meas}\prod_{\lambda_n\neq0}\lambda_n,
\end{equation*}
Grassmann 规则 $\int\dd a_r=0$、$\int\dd a_r\,a_r=1$ 表明：没有外场插入时该拓扑扇区的费米子积分为零；一个非零相关函数必须为每个零模恰好提供一个相应的场因子。对 $\nu=1$，每个左手弱二重态必须出现一次。

一代标准模型含三个颜色副本 $Q_L^a$（$a=1,2,3$）和一个 $L_L$。把颜色收缩成单态、弱同位旋指标用 $\epsilon_{ij}$ 配对，零模饱和产生的局域多费米子结构可示意写成
\begin{equation*}
 \mathfrak O_{\rm tH}^{(g)}
 \equiv
 \epsilon_{abc}
 \bigl(Q_{Lg}^{a}Q_{Lg}^{b}\bigr)
 \bigl(Q_{Lg}^{c}L_{Lg}\bigr),
\end{equation*}
其中括号表示 Lorentz 与 $SU(2)_L$ 指标按不变张量收缩。三个夸克场各带 $B=1/3$，一个轻子场带 $L=1$，所以单代、单位拓扑荷的顶点改变
\begin{equation*}
 \Delta B_{g}=3\times\frac13=1,
 \qquad
 \Delta L_{g}=1.
\end{equation*}

同一结果也可直接从反常 Ward 恒等式积分得到。对重子流和轻子流，零模计数给每代相同的拓扑密度系数，
\begin{equation*}
 \partial_\mu j_B^\mu
 =N_g\frac{g^2}{32\pi^2}
 \mathcal W_{\mu\nu}^I{}^\star\!\mathcal W^{I\mu\nu},
 \qquad
 \partial_\mu j_L^\mu
 =N_g\frac{g^2}{32\pi^2}
 \mathcal W_{\mu\nu}^I{}^\star\!\mathcal W^{I\mu\nu}.
\end{equation*}
在一个连接两个渐近真空的四维区域上积分，并使用正文对 $\nu$ 的定义，
\begin{align*}
 \Delta B
 &=\int\dd^4x\,\partial_\mu j_B^\mu
 =N_g\nu,\\
 \Delta L
 &=\int\dd^4x\,\partial_\mu j_L^\mu
 =N_g\nu.
\end{align*}
对任意整数 $\nu$，$|\nu|$ 个零模副本给出同样的线性倍数；$\nu<0$ 时手征性与粒子数变化同时反向。于是
\begin{equation*}
 \Delta(B-L)=0,
 \qquad
 \Delta(B+L)=2N_g\nu,
\end{equation*}
即正文\eqnrefs{eq:sm-bplusl-selection-dp32}。选择定则来自费米子测度在不同拓扑扇区中的零模结构，而不是规范反常未抵消。
\end{derivation}

\begin{derivation}[轴向重定义与强相互作用的 CP 相位]
把所有有色 Dirac 夸克收成味向量 $q=(u,d,s,\ldots)^T$，质量项写成
\begin{equation*}
 \Lag_m=-\bar q_LMq_R-\bar q_RM^\dagger q_L,
\end{equation*}
其中 $M$ 可以是任意非奇异复矩阵。取对角的实味空间矩阵
\begin{equation*}
 A={\rm diag}(\alpha_1,\ldots,\alpha_{N_f})
\end{equation*}
并作独立轴向重相位
\begin{equation*}
 q_L\longrightarrow \ee^{-\ii A}q_L,
 \qquad
 q_R\longrightarrow \ee^{+\ii A}q_R.
\end{equation*}
因为 $\bar q_L\to\bar q_Le^{+\ii A}$，质量矩阵在新场基底中变成
\begin{equation*}
 M' =\ee^{+\ii A}Me^{+\ii A}.
\end{equation*}
对行列式取相位，利用 $\det(\ee^{\ii A})=\ee^{\ii\Tr A}$：
\begin{align*}
 \det M'
 &=\ee^{\ii\Tr A}\det M\,\ee^{\ii\Tr A}
 =\ee^{2\ii\Tr A}\det M,\\
 \arg\det M'
 &=\arg\det M+2\Tr A
 \pmod{2\pi}.
\end{align*}
因此轴向旋转确实可以移动夸克质量矩阵的整体复相位；矢量型重相位 $q_L,q_R\to \ee^{\ii B}q_{L,R}$ 则不会产生这一相位变化。

另一方面，正文\eqnrefs{eq:qcd-singlet-axial-anomaly-dp32} 表明同一个轴向变换不是量子测度的对称。正文\eqnrefs{eq:qcd-theta-action-dp89} 已固定 $\theta_s$ 的归一化。对第 $i$ 个基本表示 Dirac 味作角度 $\alpha_i$ 的轴向旋转，Fujikawa Jacobian 给
\begin{equation*}
 \ln J_i
 =-\ii\alpha_i\frac{g_s^2}{16\pi^2}
 \int\dd^4x\,
 \mathcal F_{\mu\nu}^a{}^\star\!\mathcal F^{a\mu\nu}.
\end{equation*}
把 $J_i$ 吸收到 $\exp(\ii S/\hbar)$ 中，等价于
\begin{equation*}
 \theta_s\longrightarrow\theta_s-2\alpha_i.
\end{equation*}
全部味同时变换时便有
\begin{equation*}
 \theta_s' =\theta_s-2\Tr A.
\end{equation*}
与质量行列式的变化相加：
\begin{align*}
 \theta_s'+\arg\det M'
 &=(\theta_s-2\Tr A)
 +(\arg\det M+2\Tr A)\\
 &=\theta_s+\arg\det M.
\end{align*}
所以轴向基底选择能分别改变 $\theta_s$ 与质量矩阵的整体相位，却不能改变组合
\begin{equation*}
 {\bar\theta=\theta_s+\arg\det M}.
\end{equation*}
对标准模型六个夸克，把 $M$ 分成上、下型两个块就得到正文\eqnrefs{eq:strong-cp-theta-bar-dp11} 的 $\arg\det(M_uM_d)$。

若某个夸克严格无质量，例如某个质量本征值 $m_j=0$，则相应轴向角 $\alpha_j$ 不会在质量项中产生可观测相位：$\det M=0$ 时整体行列式相位本身不再是物理参数。此时可以选择 $\alpha_j=\theta_s/2$ 把拓扑角完全移入这个无质量场的轴向定义，从而使 $\theta_s$ 不可观测。真实 QCD 中所有夸克质量均非零，因此这一消除机制不存在，$\bar\theta$ 成为独立的强 $CP$ 参数。
\end{derivation}

\begin{derivation}[味求和中的 GIM 抵消与质量抑制]
质量本征基中的夸克带电流为
\begin{equation*}
 \Lag_{\rm CC}^{q}
 =-\frac{\hbar c\,g}{\sqrt2}
 \left[
 \bar u_i\gamma^\mu P_LV_{ij}d_j\,\mathcal W_\mu^+
 +\bar d_j\gamma^\mu P_LV_{ij}^*u_i\,\mathcal W_\mu^-
 \right].
\end{equation*}
考虑 $d_j\to d_k$、$j\neq k$ 的一圈味改变中性流。Dyson 展开中任意包含两个带电流顶角的收缩，都含味因子
\begin{equation*}
 V_{ij}V_{ik}^*=V_{ik}^*V_{ij},
\end{equation*}
因为第一个顶角把外部 $d_j$ 连接到内部 $u_i$，第二个厄米共轭顶角再把同一内部味 $u_i$ 连接到外部 $d_k$。

给定外部粒子与动量后，选一组完备、线性无关的 Lorentz--Dirac 算符基 $\{\mathcal K_a(p_{\rm ext})\}$。同一微扰阶的全部规范完成图求和后，任意内部味 $i$ 的贡献都能按这组基展开：
\begin{equation*}
 \mathcal A_{j\to k}^{(i)}
 =V_{ik}^*V_{ij}
 \sum_a\mathcal K_a(p_{\rm ext})F_a(x_i),
 \qquad
 x_i=\frac{m_{u_i}^2}{m_W^2}.
\end{equation*}
这里 $F_a$ 由实际圈积分定义。例如在维数正规化中，每个 $F_a$ 都来自形如
\begin{equation*}
 \mu_{\rm DR}^{2\epsilon}
 \int\frac{\dd^d\ell}{(2\pi)^d}
 \frac{N_a(\ell,p_{\rm ext};m_{u_i},m_W)}
 {\prod_r D_r(\ell,p_{\rm ext};m_{u_i},m_W)}
\end{equation*}
的积分及同阶规范完成项。对内部味求和并交换有限求和与算符基展开，得到
\begin{align*}
 \mathcal A_{j\to k}^{\rm FCNC}
 &=\sum_i\mathcal A_{j\to k}^{(i)}\\
 &=\sum_a\mathcal K_a(p_{\rm ext})
 \left[\sum_iV_{ik}^*V_{ij}F_a(x_i)\right].
\end{align*}
因此
\begin{equation*}
 C_a^{jk}=\sum_iV_{ik}^*V_{ij}F_a(x_i),
\end{equation*}
即正文\eqnrefs{eq:sm-gim-loop-sum-dp83}。

CKM 幺正性 $V^\dagger V=I$ 的非对角元满足
\begin{equation*}
 \sum_iV_{ik}^*V_{ij}=0,
 \qquad j\neq k.
\end{equation*}
对任意参考点 $x_*$，逐个算符系数加减同一个 $F_a(x_*)$：
\begin{align*}
 C_a^{jk}
 &=\sum_iV_{ik}^*V_{ij}
 \left[F_a(x_i)-F_a(x_*)+F_a(x_*)\right]\\
 &=\sum_iV_{ik}^*V_{ij}[F_a(x_i)-F_a(x_*)]
 +F_a(x_*)\sum_iV_{ik}^*V_{ij}\\
 &=\sum_iV_{ik}^*V_{ij}[F_a(x_i)-F_a(x_*)].
\end{align*}
这就是\eqnrefs{eq:sm-gim-subtracted-exact-dp85}。若 $x_i=x_*$ 对所有内部味都成立，则每个独立算符系数 $C_a^{jk}$ 都等于零，因而完整振幅为零。

若 $F_a$ 在连接 $x_*$ 与 $x_i$ 的区间可微，微积分基本定理还给出精确质量差表示
\begin{align*}
 F_a(x_i)-F_a(x_*)
 &=\int_{x_*}^{x_i}\dd y\,F_a'(y)\\
 &=(x_i-x_*)\int_0^1\dd t\,
 F_a'\!\left[x_*+t(x_i-x_*)\right],
\end{align*}
因此
\begin{equation*}
 C_a^{jk}
 =\sum_iV_{ik}^*V_{ij}(x_i-x_*)
 \int_0^1\dd t\,
 F_a'\!\left[x_*+t(x_i-x_*)\right].
\end{equation*}
这条等式直接表明共同的内部质量部分被消去，剩余量由质量劈裂加权。

对轻内部质量取 $x_*=0$。若 $F_a\in C^3([0,x_{\max}])$，Taylor 定理带积分余项给
\begin{equation*}
 F_a(x_i)
 =F_a(0)+x_iF_a'(0)+\frac{x_i^2}{2}F_a''(0)
 +\frac{x_i^3}{2}\int_0^1\dd t\,(1-t)^2F_a^{(3)}(tx_i).
\end{equation*}
代入并使用 $\sum_iV_{ik}^*V_{ij}=0$：
\begin{equation*}
 C_a^{jk}
 =F_a'(0)\sum_iV_{ik}^*V_{ij}x_i
 +\mathcal R_{a,jk}^{(\ge2)},
\end{equation*}
其中
\begin{align*}
 \mathcal R_{a,jk}^{(\ge2)}
 &\equiv
 \frac{F_a''(0)}2
 \sum_iV_{ik}^*V_{ij}x_i^2+R_{a,jk}^{(3)},\\
 R_{a,jk}^{(3)}
 &\equiv
 \frac12\sum_iV_{ik}^*V_{ij}x_i^3
 \int_0^1\dd t\,(1-t)^2F_a^{(3)}(tx_i).
\end{align*}
于是
\begin{align*}
 |R_{a,jk}^{(3)}|
 &\le\frac16
 \max_{0\le y\le x_{\max}}|F_a^{(3)}(y)|
 \sum_i|V_{ik}^*V_{ij}|x_i^3,\\
 |\mathcal R_{a,jk}^{(\ge2)}|
 &\le\frac{|F_a''(0)|}{2}
 \sum_i|V_{ik}^*V_{ij}|x_i^2
 +|R_{a,jk}^{(3)}|.
\end{align*}
最后用 $x_i=m_{u_i}^2/m_W^2$ 即得到正文\eqnrefs{eq:sm-gim-mass-suppression-dp83} 及其余项控制。若 $F_a$ 在零点附近非解析，或某个 $x_i$ 不小，则保留精确差值\eqnrefs{eq:sm-gim-subtracted-exact-dp85}。
\end{derivation}

\begin{derivation}[场重定义后的独立物理参数]
Yukawa 矩阵的矩阵元依赖味基选择，参数计数必须先除去这种场重定义冗余，因此从夸克扇区开始。三代时，$Y_u$ 与 $Y_d$ 各是一般复 $3\times3$ 矩阵。每个矩阵含

\begin{equation*}
 9\ \text{个复矩阵元}
 =18\ \text{个实参数},
\end{equation*}
所以两者合计有 $36$ 个实参数。

暂时令 Yukawa 耦合为零。夸克动能项在世代空间具有
\begin{equation*}
 U(3)_{Q_L}\times U(3)_{u_R}\times U(3)_{d_R}
\end{equation*}
的味基冗余。每个 $U(3)$ 有 $3^2=9$ 个连续生成元，因此三组基变换一共有 $27$ 个连续参数。不过不能直接用 $36-27$，因为一般非零的 $Y_u,Y_d$ 仍保留共同的矢量式重子数变换
\begin{equation*}
 Q_L\to \ee^{\ii\alpha}Q_L,
 \qquad
 u_R\to \ee^{\ii\alpha}u_R,
 \qquad
 d_R\to \ee^{\ii\alpha}d_R.
\end{equation*}
在 Yukawa 双线性项中，左手场经 Dirac 共轭后产生 $\ee^{-\ii\alpha}$，恰与右手场的 $\ee^{+\ii\alpha}$ 相消。因此 $U(1)_B$ 是打开 Yukawa 耦合后仍未破缺的子群，不能用来改变 $Y_u,Y_d$。真正能够消去坐标冗余的味基参数只有
\begin{equation*}
 27-1=26.
\end{equation*}
所以夸克 Yukawa 扇区留下
\begin{equation*}
 36-26=10
\end{equation*}
个物理参数。为了看清这 $10$ 个参数为什么恰好表现为“六个质量、三个角、一个相位”，还要把上述商空间计数落实到质量本征基。对两个一般复矩阵分别作奇异值分解：
\begin{equation*}
 Y_u=U_{uL}\,y_u\,U_{uR}^\dagger,
 \qquad
 Y_d=U_{dL}\,y_d\,U_{dR}^\dagger,
\end{equation*}
其中
\begin{equation*}
 y_u=\operatorname{diag}(y_u,y_c,y_t),
 \qquad
 y_d=\operatorname{diag}(y_d,y_s,y_b)
\end{equation*}
可取为非负实对角矩阵。右手幺正矩阵 $U_{uR},U_{dR}$ 把两个质量矩阵对角化；带电弱流只含左手双重态，因此两个左手旋转之间的相对矩阵保留下来：
\begin{equation*}
 V_{\rm CKM}=U_{uL}^\dagger U_{dL}.
\end{equation*}
六个奇异值与六个夸克质量一一对应，剩下的问题就是一个一般 $3\times3$ 幺正矩阵究竟含多少可观测参数。

一般 $U(3)$ 有 $9$ 个实参数。可把它们组织成三个实混合角与六个相位。质量项已经对角以后，六个 Dirac 夸克仍可分别作共同左右手重相位
\begin{equation*}
 u_i\to \ee^{\ii\alpha_i}u_i,
 \qquad
 d_j\to \ee^{\ii\beta_j}d_j,
\end{equation*}
而不改变任何质量本征值。带电流
$\bar u_{Li}\gamma^\mu V_{ij}d_{Lj}$ 因而只把 CKM 元素变换为
\begin{equation*}
 V_{ij}\longrightarrow
 \ee^{-\ii\alpha_i}V_{ij}\ee^{+\ii\beta_j}.
\end{equation*}
六个相位 $\{\alpha_i,\beta_j\}$ 中，若全部同时增加同一个 $\omega$，则
\begin{equation*}
 \ee^{-\ii(\alpha_i+\omega)}V_{ij}
 \ee^{+\ii(\beta_j+\omega)}
 =\ee^{-\ii\alpha_i}V_{ij}\ee^{+\ii\beta_j},
\end{equation*}
所以共同的重子数相位并不能改变 $V_{\rm CKM}$。真正可用于消去 CKM 相位的只有 $6-1=5$ 个独立重相位。因此六个原始相位中只剩
\begin{equation*}
 6-5=1
\end{equation*}
个物理 $CP$ 相位，而三个实混合角保留。于是夸克 Yukawa 扇区的 $10$ 个参数被显式分解为
\begin{equation*}
 {6\ \text{个夸克质量}
 +3\ \text{个 CKM 混合角}
 +1\ \text{个 CKM 相位}.}
\end{equation*}

带电轻子扇区只有一个一般复 $3\times3$ 矩阵 $Y_e$，起始共有 $18$ 个实参数。令 Yukawa 耦合为零时，动能项具有
\begin{equation*}
 U(3)_{L_L}\times U(3)_{e_R}
\end{equation*}
的 $18$ 参数味基冗余。严格的最小标准模型中没有右手中微子，中微子保持无质量；把 $Y_e$ 对角化以后，三个轻子家族仍可分别做保持动能项和质量项不变的共同矢量式重相位。因此未破缺子群是 $U(1)^3$，维数为 $3$。可用于消去 $Y_e$ 坐标冗余的基变换参数数目为
\begin{equation*}
 18-3=15,
\end{equation*}
于是只剩
\begin{equation*}
 18-15=3
\end{equation*}
个带电轻子质量；在这个严格最小场内容中没有物理的 PMNS 混合角与相位。

所以全部 Yukawa 与味结构一共贡献
\begin{equation*}
 10+3=13
\end{equation*}
个物理参数。再加入三种规范耦合 $g_s,g,g'$、Higgs 势中的两个参数 $\mu_H^2,\lambda_H$，以及一个基不变的强相互作用 $CP$ 参数 $\bar\theta$，总数为
\begin{equation*}
 13+3+2+1=19.
\end{equation*}
于是得到\eqnrefs{eq:minimal-sm-parameter-count-dp11}。19 个连续参数就是耦合参数空间对场重定义轨道取商后的维数；粒子质量、混合参数和规范耦合只是这个商空间的一组物理坐标。
\end{derivation}

\begin{derivation}[质量基中的 Higgs 费米子耦合]
对任意一种带电费米子 $f=u,d,e$，电弱破缺后正文 Yukawa 项都具有同一矩阵结构
\begin{equation*}
 \Lag_{Y,f}
 =-\frac{v_E}{\sqrt2}
 \left(\bar f_LY_ff_R+\bar f_RY_f^\dagger f_L\right)
 -\frac{\sqrt{\hbar c}\,h}{\sqrt2}
 \left(\bar f_LY_ff_R+\bar f_RY_f^\dagger f_L\right).
\end{equation*}
这里 $Y_f$ 是世代空间中的任意复 $N_g\times N_g$ 矩阵；它一般既不厄米也不正规，因此不能假定存在单个幺正矩阵 $U$ 使 $U^\dagger Y_fU$ 对角。

从两个厄米半正定矩阵
\begin{equation*}
 H_L\equiv Y_fY_f^\dagger,
 \qquad
 H_R\equiv Y_f^\dagger Y_f
\end{equation*}
出发。谱定理保证存在幺正矩阵 $U_{fL}$ 使
\begin{equation*}
 U_{fL}^\dagger H_LU_{fL}
 =\operatorname{diag}(y_1^2,\ldots,y_{N_g}^2),
 \qquad y_i\ge0.
\end{equation*}
设 $u_{Li}$ 是 $U_{fL}$ 的第 $i$ 个列向量，即
\begin{equation*}
 H_Lu_{Li}=y_i^2u_{Li}.
\end{equation*}
对任意 $y_i>0$，定义右奇异向量
\begin{equation*}
 u_{Ri}\equiv\frac1{y_i}Y_f^\dagger u_{Li}.
\end{equation*}
它们的内积为
\begin{align*}
 u_{Ri}^\dagger u_{Rj}
 &=\frac1{y_iy_j}
 u_{Li}^\dagger Y_fY_f^\dagger u_{Lj}\\
 &=\frac{y_j^2}{y_iy_j}u_{Li}^\dagger u_{Lj}
 =\delta_{ij},
\end{align*}
并且
\begin{align*}
 Y_fu_{Ri}
 &=\frac1{y_i}Y_fY_f^\dagger u_{Li}
 =y_i u_{Li},\\
 Y_f^\dagger u_{Li}
 &=y_i u_{Ri}.
\end{align*}
若存在 $y_i=0$，在 $\ker Y_f$ 与 $\ker Y_f^\dagger$ 中分别补齐正交基即可。把全部 $u_{Ri}$ 作为列向量组成幺正矩阵 $U_{fR}$，便有逐矩阵元关系
\begin{align*}
 (U_{fL}^\dagger Y_fU_{fR})_{ij}
 &=u_{Li}^\dagger Y_fu_{Rj}\\
 &=y_j u_{Li}^\dagger u_{Lj}
 =y_j\delta_{ij}.
\end{align*}
因此
\begin{equation*}
 {
 U_{fL}^\dagger Y_fU_{fR}
 =\operatorname{diag}(y_1,\ldots,y_{N_g})}
\end{equation*}
就是正文\eqnrefs{eq:yukawa-svd-dp11}。奇异值按定义非负，因此由 $m_f=v_E y_f/(\sqrt2 c)$ 得到的质量本征值可统一取非负实数。

定义新的左右手场
\begin{equation*}
 f_L'=U_{fL}^\dagger f_L,
 \qquad
 f_R'=U_{fR}^\dagger f_R.
\end{equation*}
因为 $U_{fL},U_{fR}$ 只作用在世代空间而且为常数矩阵，动能保持不变：
\begin{align*}
 \bar f_L\ii\hbar c\gamma^\mu D_\mu f_L
 &=\bar f_L'U_{fL}^\dagger
 \ii\hbar c\gamma^\mu D_\mu U_{fL}f_L'\\
 &=\bar f_L'\ii\hbar c\gamma^\mu D_\mu f_L',
\end{align*}
右手项同理。这里使用的是规范相互作用在同一电荷扇区的世代空间中正比于单位矩阵，因此 $D_\mu$ 与 $U_{fL,R}$ 对易。

质量项变为
\begin{align*}
 -\frac{v_E}{\sqrt2}\bar f_LY_ff_R+\mathrm{H.c.}
 &=-\frac{v_E}{\sqrt2}
 \bar f_L'U_{fL}^\dagger Y_fU_{fR}f_R'
 +\mathrm{H.c.}\\
 &=-\sum_i\frac{y_iv_E}{\sqrt2}
 \left(\bar f_{Li}'f_{Ri}'+\bar f_{Ri}'f_{Li}'\right)\\
 &=-\sum_i m_{fi}c^2\,\bar f_i'f_i',
\end{align*}
其中
\begin{equation*}
 {m_{fi}c^2=\frac{y_iv_E}{\sqrt2}},
\end{equation*}
即正文\eqnrefs{eq:fermion-masses-yukawa-dp11}。

关键点是 Higgs 涨落项含有\emph{完全相同的}矩阵 $Y_f$，因此不需要第二次对角化：
\begin{align*}
 \Lag_{hff}
 &=-\frac{\sqrt{\hbar c}\,h}{\sqrt2}
 \bar f_L'U_{fL}^\dagger Y_fU_{fR}f_R'
 +\mathrm{H.c.}\\
 &=-\sum_i
 \frac{y_i\sqrt{\hbar c}}{\sqrt2}
 h\left(\bar f_{Li}'f_{Ri}'+\bar f_{Ri}'f_{Li}'\right)\\
 &=-\sum_i
 \frac{m_{fi}c^2}{v_E}\sqrt{\hbar c}\,
 h\,\bar f_i'f_i'.
\end{align*}
这就是正文\eqnrefs{eq:higgs-fermion-coupling-dp11}。所以树级 Higgs 中性耦合在质量基中自动对角，且耦合强度与费米子质量线性成正比。

若某组奇异值严格简并，则对应简并子空间内仍允许额外幺正旋转而不改变质量矩阵；因此 CKM/PMNS 的独立参数计数需要单独注明“质量谱非简并”这一条件。质量本征值的定义和 Higgs 中性耦合的对角性本身并不依赖这一假设。
\end{derivation}

\begin{derivation}[一代费米子的规范反常抵消]
规范反常的群因子来自三角费米子圈。对一个左手 Weyl 费米子，三个外部规范玻色子与费米子环相连时，两种相反顶角次序的群因子分别为
\begin{equation*}
 \Tr_R(T^aT^bT^c),
 \qquad
 \Tr_R(T^aT^cT^b).
\end{equation*}
三角图的宇称奇部分在交换后两个外腿时是对称的，因此两种次序相加后只保留
\begin{equation*}
 \Tr_R\!\left[T^a\{T^b,T^c\}\right]
 =\frac12d_R^{abc}.
\end{equation*}
动量积分与正则化决定共同的数值前因子，而一个给定费米子谱是否无规范反常只取决于全部左手 Weyl 场的这些群因子之和是否为零。若其中一个或多个外腿属于 $U(1)_Y$，相应生成元就是超荷乘单位矩阵，于是四类局域反常系数分别化为
\begin{align*}
 [SU(N)]^2U(1)_Y &: \sum_f Y_fT(R_f),\\
 [U(1)_Y]^3 &: \sum_f d(R_f)Y_f^3,\\
 {\rm grav}^2U(1)_Y &: \sum_f d(R_f)Y_f,\\
 [SU(N)]^3 &: \sum_f d_{R_f}^{abc}.
\end{align*}
这里 $d(R_f)$ 表示其余内部指标的多重度。规范理论要求这些群论系数在全部左手 Weyl 场上求和后同时为零。

现在把一代标准模型全部改写成左手 Weyl 场
\begin{equation*}
 Q_L:(\mathbf3,\mathbf2)_{1/6},\quad
 u_R^c:(\bar{\mathbf3},\mathbf1)_{-2/3},\quad
 d_R^c:(\bar{\mathbf3},\mathbf1)_{1/3},\quad
 L_L:(\mathbf1,\mathbf2)_{-1/2},\quad
 e_R^c:(\mathbf1,\mathbf1)_{1}.
\end{equation*}
把共同的三角圈动量积分和规范耦合常数提出以后，定义纯群论反常系数
\begin{align*}
 \mathfrak A_{33Y}
 &\equiv\sum_{\text{左手 Weyl }f}d_2(f)Y_fT_3(R_f),\\
 \mathfrak A_{22Y}
 &\equiv\sum_f d_3(f)Y_fT_2(R_f),\\
 \mathfrak A_{YYY}
 &\equiv\sum_f d_3(f)d_2(f)Y_f^3,\\
 \mathfrak A_{{\rm grav}^2Y}
 &\equiv\sum_f d_3(f)d_2(f)Y_f.
\end{align*}
其中 $d_3,d_2$ 分别是颜色和弱表示维数。利用 $T(\mathbf3)=T(\bar{\mathbf3})=T(\mathbf2)=1/2$，一代标准模型逐项给
\begin{align*}
 \mathfrak A_{33Y}
 &=2\!\left(\frac16\right)\frac12
 +\left(-\frac23\right)\frac12
 +\left(\frac13\right)\frac12=0,\\
 \mathfrak A_{22Y}
 &=3\!\left(\frac16\right)\frac12
 +\left(-\frac12\right)\frac12=0,\\
 \mathfrak A_{YYY}
 &=6\left(\frac16\right)^3
 +3\left(-\frac23\right)^3
 +3\left(\frac13\right)^3
 +2\left(-\frac12\right)^3+1^3=0,\\
 \mathfrak A_{{\rm grav}^2Y}
 &=6\left(\frac16\right)
 +3\left(-\frac23\right)
 +3\left(\frac13\right)
 +2\left(-\frac12\right)+1=0.
\end{align*}
这些等式就是四类反常的群论系数本身为零，而不是只说明它们“与某个零表达式成正比”。
对 $[SU(3)_c]^3$，基本表示与反基本表示的 $d_R^{abc}$ 符号相反，而 $Q_L$ 有两个弱分量，因此
\begin{equation*}
 2d_{\mathbf3}^{abc}+d_{\bar{\mathbf3}}^{abc}+d_{\bar{\mathbf3}}^{abc}
 =(2-1-1)d_{\mathbf3}^{abc}=0.
\end{equation*}
$SU(2)$ 基本表示是伪实表示，局域三次群不变量本身为零。因此一代费米子谱的全部局域规范反常系数均消失，即得到正文\eqnrefs{eq:sm-gauge-anomaly-cancellation-dp32}。正文另行陈述的 Witten 模二条件属于大规范变换的全局拓扑定理，不由这个微扰三角图计算推出。
\end{derivation}
